\chapter{Principal angles, projectors, and overlap between subspaces}

\section{The one-dimensional case}

For two unit vectors \(\mathbf u_1\) and \(\mathbf u_2\), the acute angle is
\begin{equation}
\psi=\arccos\left(|\mathbf u_1^\top\mathbf u_2|\right).
\end{equation}
The absolute value appears because a line has no inherent sign: \(\mathbf u\) and \(-\mathbf u\) span the same one-dimensional subspace.

\inwords The inner product is the length of one unit vector retained when it is projected onto the other. Perfect alignment gives one; orthogonality gives zero.

\section{Principal angles for \(k\)-dimensional channels}

Let \(\Vv_1,\Vv_2\in\R^{n\times k}\) contain orthonormal bases:
\begin{equation}
\Vv_i^\top\Vv_i=\Ii_k.
\end{equation}
Form the cross-basis matrix
\begin{equation}
\mathbf C=\Vv_1^\top\Vv_2.
\end{equation}
Take its singular value decomposition:
\begin{equation}
\mathbf C=\mathbf L\diag(\sigma_1,\ldots,\sigma_k)\mathbf R^\top,
\qquad
1\ge\sigma_1\ge\cdots\ge\sigma_k\ge0.
\end{equation}
The principal angles are defined by
\begin{equation}
\boxed{\cos\psi_j=\sigma_j.}
\label{eq:principalAngles}
\end{equation}

\begin{mathteach}
The matrices \(\mathbf L\) and \(\mathbf R\) merely choose the paired directions within the two subspaces that align best, second-best, and so on. The singular values are invariant to arbitrary rotations of either basis and therefore belong to the subspaces themselves.
\end{mathteach}

\section{Projector form of chordal distance}

Let
\begin{equation}
\Pp_i=\Vv_i\Vv_i^\top.
\end{equation}
First compute
\begin{align}
\tr(\Pp_1\Pp_2)
&=\tr(\Vv_1\Vv_1^\top\Vv_2\Vv_2^\top)\\
&=\tr(\Vv_1^\top\Vv_2\Vv_2^\top\Vv_1)\\
&=\tr(\mathbf C\mathbf C^\top)\\
&=\normF{\mathbf C}^2\\
&=\sum_{j=1}^k\sigma_j^2\\
&=\sum_{j=1}^k\cos^2\psi_j.
\label{eq:traceProjectorOverlap}
\end{align}
The second line uses cyclic invariance of trace. The fifth uses the fact that the squared Frobenius norm is the sum of squared singular values.

Now expand the projector difference:
\begin{align}
\normF{\Pp_1-\Pp_2}^2
&=\tr\left[(\Pp_1-\Pp_2)^2\right]\\
&=\tr(\Pp_1^2)+\tr(\Pp_2^2)-2\tr(\Pp_1\Pp_2).
\end{align}
Because projectors satisfy \(\Pp_i^2=\Pp_i\) and \(\tr(\Pp_i)=k\),
\begin{align}
\normF{\Pp_1-\Pp_2}^2
&=2k-2\tr(\Pp_1\Pp_2)\\
&=2\sum_{j=1}^k(1-\cos^2\psi_j)\\
&=2\sum_{j=1}^k\sin^2\psi_j.
\end{align}
Therefore the squared chordal distance is
\begin{equation}
\boxed{
d_C^2(\mathcal C_1,\mathcal C_2)
=\half\normF{\Pp_1-\Pp_2}^2
=\sum_{j=1}^k\sin^2\psi_j.}
\label{eq:chordalIdentity}
\end{equation}

\inwords Chordal distance is the total squared amount by which the paired channel directions fail to project onto one another.

\begin{geombox}
The projector map embeds the curved Grassmann manifold into the flat Euclidean space of symmetric matrices. Chordal distance is ordinary Euclidean distance between the embedded projectors, scaled by one half.
\end{geombox}

\section{Volume overlap}

From the SVD,
\begin{align}
\det(\Vv_1^\top\Vv_2)^2
&=\det(\mathbf L)^2
\det[\diag(\sigma_j)]^2
\det(\mathbf R)^2\\
&=\prod_{j=1}^k\sigma_j^2\\
&=\prod_{j=1}^k\cos^2\psi_j,
\end{align}
because orthogonal matrices have determinant \(\pm1\).

Define the bounded channel-overlap deficit
\begin{equation}
\boxed{
I_V=1-\det(\Vv_1^\top\Vv_2)^2
=1-\prod_{j=1}^k\cos^2\psi_j.}
\label{eq:volumeDeficit}
\end{equation}

\inwords The product \(\prod_j\cos\psi_j\) is the factor by which projecting a unit \(k\)-dimensional volume element from one channel into the other shrinks that volume. Squaring removes orientation and gives a convenient overlap measure.

A related additive quantity is the log-overlap loss
\begin{equation}
L_V=-\log(1-I_V)
=-2\sum_{j=1}^k\log(\cos\psi_j).
\label{eq:logOverlap}
\end{equation}
It is additive across principal directions, but it is still a geometric descriptor rather than a calibrated fitness load.

\section{Small-angle expansion}

For small \(\psi\),
\begin{align}
\sin^2\psi
&=\psi^2-\frac{\psi^4}{3}+O(\psi^6),\\
\cos^2\psi
&=1-\psi^2+\frac{\psi^4}{3}+O(\psi^6).
\end{align}
Therefore
\begin{equation}
d_C^2
=\sum_j\psi_j^2-\frac13\sum_j\psi_j^4+O(\normtwo{\psiv}^6).
\label{eq:chordalExpansion}
\end{equation}
For the volume deficit,
\begin{align}
\prod_j\cos^2\psi_j
&=\prod_j\left(1-\psi_j^2+\frac{\psi_j^4}{3}+O(\psi_j^6)\right)\\
&=1-\sum_j\psi_j^2
+\frac13\sum_j\psi_j^4
+\sum_{i<j}\psi_i^2\psi_j^2
+O(\normtwo{\psiv}^6),
\end{align}
so
\begin{equation}
I_V
=\sum_j\psi_j^2
-\frac13\sum_j\psi_j^4
-\sum_{i<j}\psi_i^2\psi_j^2
+O(\normtwo{\psiv}^6).
\label{eq:volumeExpansion}
\end{equation}
Thus
\begin{equation}
\boxed{I_V=d_C^2+O(\normtwo{\psiv}^4)
=\sum_j\psi_j^2+O(\normtwo{\psiv}^4).}
\label{eq:overlapLocalEquivalence}
\end{equation}

\begin{connectionbox}
The quadratic snowball does not depend on choosing the determinant deficit rather than chordal distance. All smooth overlap deficits with a nonzero quadratic term agree locally up to scaling. The choice matters at larger angles and for interpretation, but not for the early exponent.
\end{connectionbox}

\begin{figure}[H]
\centering
\includegraphics[width=0.96\textwidth]{figures/pedagogy_principal_angles.pdf}
\caption{Three equivalent geometric readings of channel alignment: angle, projection length, and projected \(k\)-volume.}
\label{fig:principalAnglePedagogy}
\end{figure}
